\begin{pastexam}[label={px:gw3-103-ds-05}]{關務三等資訊處理 103 年 第五題}
（申論；20 分）

\begin{itemize}[leftmargin=2.2em]
\item[㈠] 陣列（array）、鏈結串列（linked list）是常見的線性資料結構（linear data structure）。請列舉兩種非線性的
資料結構（non-linear data structure）。（8 分）
\item[㈡] 在巨量資料（big data）分析中，何謂結構化資料（structured data）？請舉例說明（6 分）；何謂非結構化資料
（unstructured data）？請舉例說明（6 分）。
\end{itemize}
\end{pastexam}

\begin{pastsoln}
答案：㈠ 樹（tree，例如二元樹、堆積）和圖（graph）。 ㈡ 結構化資料有固定欄位、可以放進表格，例如資料庫裡的報關單
紀錄；非結構化資料沒有預先定義的欄位，例如貨櫃 X 光影像、電子郵件內文。

\textbf{㈠} 線性資料結構的元素排成一列，除了頭尾，每個元素恰有一個前一個和一個後一個（陣列、鏈結串列、堆疊、佇列）。
非線性資料結構的元素之間不是一對一的前後關係：
\begin{enumerate}
\item \textbf{樹}（tree）：階層式的一對多關係，每個節點有一個父節點（根除外），可以有多個子節點。 例如二元樹、二元
搜尋樹、堆積、B 樹，用來表示目錄結構、組織圖、運算式（資結書第~7~章）。
\item \textbf{圖}（graph）：多對多的關係，頂點之間用邊任意相連，可以有環。 例如道路網、社群網路、網頁連結
（資結書第~21~章）。
\end{enumerate}
（樹其實是沒有環的連通圖的一種。 雜湊表、集合等也常被歸為非線性，答兩種時以樹和圖最標準。）

\textbf{㈡} 這一小題屬資料管理的觀念，不是資料結構本身的內容，簡要回答即可。
\begin{itemize}
\item \textbf{結構化資料}：有事先定義好的格式（schema），每筆資料都有相同的欄位、每個欄位有固定的型別，可以整齊地
放進列與行的表格，用 SQL 這類查詢語言直接篩選、排序、彙總。 例如：海關的報關資料表，每筆有報單號碼、日期、
貨品分類號列、數量、完稅價格等欄位；銀行的交易紀錄；學生成績表。
\item \textbf{非結構化資料}：沒有固定的欄位格式，無法直接放進表格，要先經過處理（文字探勘、影像辨識、語音辨識等）
才能抽出可分析的資訊。 例如：貨櫃的 X 光掃描影像、監視器影片、電子郵件和社群貼文的文字、客服通話錄音、掃描的
PDF 文件。 巨量資料中非結構化資料占大多數，處理上也比較困難。
\end{itemize}
介於兩者之間的叫\textbf{半結構化資料}：沒有固定的表格結構，但有標記可以辨識欄位，例如 JSON、XML 檔、網頁的 HTML。
\end{pastsoln}
